\documentclass[12pt]{article}
\usepackage{hyperref}

\usepackage{graphicx}

\usepackage{mathtools}
\DeclarePairedDelimiter\ceil{\lceil}{\rceil}
\DeclarePairedDelimiter\floor{\lfloor}{\rfloor}

\usepackage{cancel}

\usepackage{amsmath}
\usepackage{amsfonts}
\usepackage{amssymb}

\usepackage{listings}
\lstset{language=C++,
                basicstyle=\ttfamily,
                keywordstyle=\color{blue}\ttfamily,
                stringstyle=\color{red}\ttfamily,
                commentstyle=\color{green}\ttfamily,
                morecomment=[l][\color{magenta}]{\#}
}

\usepackage{breqn}
\usepackage[]{algorithm2e}
\usepackage{physics}

\newcommand\fnurl[2]{%
  \href{#2}{#1}\footnote{\url{#2}}%
}

\newcommand{\Four}{\mathcal{F}}
\renewcommand{\(}{\left(}
\renewcommand{\)}{\right)}
%% \renewcommand{\{}{\left{}
%% \renewcommand{\}}{\right}}


\title{MPI with rocFFT}
\author{rocFFT}



\def\no#1{\nobreak{#1}} % stop equation breaking in dmath

\begin{document}
\maketitle

\section{Introduction}

These are notes on distritubed implementations for rocFFT.

\section{Decompositions}

\url{https://www.osti.gov/servlets/purl/1483229}

Given data of size $N^3$ with $P$ mpi procs, the brick layout has
each process having
\begin{dmath}
  \label{brick3d}
  \left(\frac{N}{P^{1/3}}, \frac{N}{P^{1/3}}, \frac{N}{P^{1/3}} \right)
\end{dmath}
and this must be transformed into either a pencil decomposition
\begin{dmath}
  \label{pencil3d}
  \left(N, \frac{N}{\sqrt{P}}, \frac{N}{\sqrt{P}} \right)
\end{dmath}
or a slab decomposition,
\begin{dmath}
  \label{slab3d}
  \left(N, N, \frac{N}{P} \right)
\end{dmath}

\subsection{Brick-Pencil}
Going from distribution~\eqref{brick3d} to~\eqref{pencil3d}, each of
the $P$ procs need to receive $P^{1/3}-1$ messages of dimensions 
\begin{dmath}
  \left(\frac{N}{P^{1/3}}, \frac{N}{\sqrt{P}}, \frac{N}{\sqrt{P}} \right)
\end{dmath}
for message size of $\frac{N^3}{P^{4/3}}$.  There are therefore
$P(P^{1/3}-1)$ messages, for a total communication size of
\begin{dmath}
  \frac{N^3}{P^{4/3}} P(P^{1/3}-1) \approx N^3.
\end{dmath}

\subsection{Brick-Slab}
Going from distribution~\eqref{brick3d} to~\eqref{slab3d}, each of the
$P$ procs needs to receive $P^{2/3}-1$ messages of dimension
\begin{dmath}
  \left(\frac{N}{P}, \frac{N}{P^{1/3}}, \frac{N}{P^{1/3}} \right)
\end{dmath}
which has message size $\frac{N^2}{P^{5/3}}$.  The total communication
  size is
\begin{dmath}
  P\(P^{2/3}-1\) \frac{N^2}{P^{5/3}} \approx N^3.
\end{dmath}

\subsection{Pencil-Pencil}
A pencil-to-pencil transpose requires all $P$ procs to send
$\sqrt{P}-1$ messages of size
\begin{dmath}
  \left(\frac{N}{\sqrt{P}}, \frac{N}{\sqrt{P}}, \frac{N}{\sqrt{P}} \right)
\end{dmath}
which has message size $\frac{N^3}{P^{3/2}}$.  The total communication
size is
\begin{dmath}
  P \(\sqrt{P}-1\) \frac{N^3}{P^{3/2}} \approx N^3.
\end{dmath}

\subsection{Pencil-Slab}
All $P$ procs send $P-1$ messages of size
\begin{dmath}
  \left(\frac{N}{P}, \frac{N}{\sqrt{P}}, \frac{N}{\sqrt{P}} \right)
\end{dmath}
for a message size of $N^3/P^2$.  The total communication size is
\begin{dmath}
  P \(P-1\) \frac{N^3}{P^2} \approx N^3.
\end{dmath}

\subsection{Slab-Slab}
All $P$ procs send $P-1$ messages of size
\begin{dmath}
  \left(\frac{N}{P}, \frac{N}{P}, N \right)
\end{dmath}
for a message size of $N^3/P^2$. The total communication size is
\begin{dmath}
  P \(P-1\) \frac{N^3}{P^2} \approx N^3.
\end{dmath}

\subsection{Transform summary}

\begin{centering}
  \begin{table}[htbp]
    \begin{tabular}{|l|l|l|}
      \hline
      Type          & Message size  & Messages per proc \\
      \hline
      Brick-pencil  & $N^3/P^{5/3}$ & $P^{1/3}$          \\
      Brick-slab    & $N^3/P^{5/3}$ & $P^{1/3}$          \\
      Pencil-pencil & $N^3/P^{3/2}$ & $\sqrt{P}$        \\
      Pencil-slab   & $N^3/P^2$    & $P$               \\
      Slab-slab     & $N^3/P^2$    & $P$               \\
      \hline
    \end{tabular}
  \end{table}
\end{centering}

\end{document}
